\chapter{Revisão Bibliográfica}\label{cap:revisao}
% ============================================================================

Três corpos de literatura sustentam o trabalho e organizam a revisão a seguir: a hipótese de eficiência de mercado e suas implicações para a previsibilidade de preços, o uso de NLP em finanças, e os protocolos de avaliação em \textit{machine learning} para séries temporais financeiras.

\section{Hipótese de eficiência de mercado}

A hipótese de eficiência de mercado (HEM), formalizada por \citeonline{fama1970}, postula que os preços dos ativos refletem toda a informação disponível. Em sua forma semiforte, a HEM implica que informações públicas --- incluindo notícias financeiras --- já estão incorporadas nos preços no momento de sua publicação, tornando impossível obter retornos sistematicamente superiores com base nessas informações. Trabalhos empíricos subsequentes, contudo, documentaram anomalias que sugerem previsibilidade parcial. \citeonline{lo2004} propôs a Hipótese de Mercados Adaptativos, na qual a eficiência varia ao longo do tempo conforme as condições de mercado e o comportamento dos agentes.

Este trabalho parte dessa tensão como ponto de partida empírico: se os preços já incorporam o sentimento das notícias financeiras, um modelo treinado sobre esse sentimento não deveria superar um \textit{baseline} autoregressivo de preço (\textit{baseline} autoregressivo: modelo que usa exclusivamente dados históricos de preço como \textit{features}); se não incorporam completamente, há espaço mensurável para predição. O Capítulo~\ref{cap:sentimento} testa essa hipótese para ITUB4, PETR4 e VALE3 sob protocolo de janela única, obtendo o resultado dessa avaliação inicial --- aparentemente favorável ao segundo cenário. O Capítulo~\ref{cap:investigacao} submete esse resultado a validação cruzada \textit{expanding-window} com múltiplas sementes e testes de significância formal, revelando que o ganho sobre o \textit{baseline} autoregressivo não é estatisticamente distinguível do acaso ($p = 0{,}49$, IC 95\% $[-0{,}012;\ +0{,}018]$), revertendo a conclusão qualitativa de positiva para nula.

\section{NLP aplicado a finanças}

A aplicação de processamento de linguagem natural à predição de preços pode ser organizada em três gerações de representação textual \cite{kellyxiu2023}. A primeira utilizou léxicos de sentimento construídos manualmente, como o Loughran-McDonald \cite{loughran2011}, atribuindo polaridade a palavras individuais com base em dicionários específicos do domínio. A segunda introduziu \textit{embeddings} densos baseados em redes neurais, como Word2Vec \cite{mikolov2013} e GloVe \cite{pennington2014}, que capturam relações semânticas distribuídas mas não são específicos ao domínio financeiro. A terceira geração, dominante desde o início da década de 2020, utiliza modelos de linguagem pré-treinados e \textit{fine-tuned} para o domínio, como o FinBERT \cite{araci2019} e suas variantes regionais, incluindo o FinBERT-PT-BR para português brasileiro \cite{souza2023}.

Convém definir o que é um \textit{embedding}, conceito central à segunda e à terceira gerações. Um \textit{embedding} é a representação de uma unidade textual (palavra, sentença ou documento) como um vetor denso de números reais em um espaço de dimensão fixa, no qual a proximidade geométrica entre vetores reflete proximidade semântica entre os textos correspondentes. O Word2Vec aprende esses vetores prevendo palavras a partir de seu contexto, de modo que palavras que ocorrem em contextos semelhantes recebem vetores próximos, e o GloVe obtém efeito análogo fatorando a matriz global de coocorrência de palavras. Tais \textit{embeddings} clássicos atribuem um único vetor a cada palavra, independentemente da frase em que aparece. Modelos contextuais mais recentes, baseados na arquitetura Transformer, produzem vetores que variam conforme o contexto e podem ser agregados para representar sentenças ou artigos inteiros; é dessa classe o modelo de \textit{embeddings} genérico de 1.024 dimensões (Ollama \texttt{qwen3-embedding:4b}) empregado no Capítulo~\ref{cap:pipeline}.

É importante notar que resultados nulos ou negativos são sub-representados na literatura de NLP financeiro devido ao viés de publicação \cite{shah2019}. \citeonline{shah2019} documentaram que muitos estudos de \textit{sentiment analysis} para previsão de ações não replicam quando avaliados fora da amostra original. \citeonline{kellyxiu2023} mostraram que a maior parte do poder preditivo de variáveis textuais desaparece quando as variáveis textuais são combinadas com \textit{features} de preço em modelos com controle explícito de dimensionalidade --- resultado alinhado com os achados deste trabalho.

\citeonline{xu2018} demonstraram que a combinação de texto de \textit{tweets} com séries históricas de preço melhora a predição de direção em ações do S\&P~500, utilizando um modelo de atenção hierárquico; seu modelo StockNet atinge acurácia de 58,2\% e \textit{Matthews correlation coefficient} (MCC) de 0,081 sobre 88 ativos norte-americanos de alto volume de negociação. \citeonline{fischer2018} aplicaram LSTMs (\textit{Long Short-Term Memory}, proposta originalmente por \citeonline{hochreiter1997}) à predição de retornos diários de ações do S\&P~500, reportando retorno diário de 0,46\% e índice de Sharpe de 5,8 (antes de custos de transação), superando o \textit{buy-and-hold}; a acurácia direcional, porém, situa-se apenas marginalmente acima de 50\%. \citeonline{ding2015} combinaram \textit{embeddings} de eventos extraídos de notícias com redes convolucionais, alcançando acurácia direcional de 64,2\% na predição do índice S\&P~500 e 65,5\% na predição de ações individuais. \citeonline{hu2018} propuseram um arcabouço de atenção hierárquica para tendências de ações chinesas orientado a notícias, reportando retorno anualizado de 61,1\% para a carteira das 40 melhores ações (a acurácia de 47,8\% não é diretamente comparável aos demais valores por se tratar de uma tarefa de três classes). No contexto brasileiro, a literatura é majoritariamente exploratória: \citeonline{cavalcante2016} apresentam um levantamento abrangente de métodos de inteligência computacional aplicados a mercados financeiros, sem reportar uma métrica preditiva única.
Nenhum desses estudos reporta análise de variância de semente ou validação cruzada com múltiplas janelas temporais --- uma prática cuja ausência, como este trabalho demonstra adiante, pode reverter as conclusões qualitativas dos resultados.

\section{Arquiteturas de modelos}

Este trabalho utiliza quatro famílias de modelos, selecionadas por sua representatividade na literatura de ML financeiro. As três primeiras são redes neurais que processam a sequência dos últimos 30 dias; a quarta é um modelo de árvores sobre \textit{features} tabulares.

O LSTM (\textit{Long Short-Term Memory}) \cite{hochreiter1997} é uma rede neural recorrente que mantém um estado de memória ao longo da sequência, controlado por \textit{gates} que decidem quais informações reter ou descartar a cada passo de tempo --- mecanismo que mitiga o problema do gradiente que desaparece em sequências longas. O BiLSTM (LSTM bidirecional) processa a sequência em ambos os sentidos temporais e concatena os dois estados, sendo amplamente utilizado em predição de preços \cite{fischer2018}; é empregado aqui nas variantes de 2 camadas/128 unidades e 1 camada/32 unidades. O Transformer \cite{vaswani2017} dispensa a recorrência: seu mecanismo de atenção \textit{multi-head} avalia diretamente a relevância de cada posição da sequência em relação a todas as demais, capturando dependências entre instantes arbitrários sem viés de localidade; usa-se aqui um \textit{encoder} de 2 camadas, 4 cabeças de atenção e $d_\text{model}=64$. A TCN (\textit{Temporal Convolutional Network}) \cite{bai2018} aplica convoluções causais dilatadas --- cada saída depende apenas de instantes passados, e as dilatações expandem o campo receptivo exponencialmente com a profundidade ---, capturando padrões temporais locais com a configuração [32,32], \textit{kernel} 3 e dilatações [1,2]. Por fim, o XGBoost \cite{chen2016} é um método de \textit{gradient boosting}: constrói um conjunto de árvores de decisão de forma sequencial, em que cada nova árvore corrige os erros residuais das anteriores. Com 300 árvores, profundidade 4 e \textit{learning rate} 0,05, é o padrão de fato para dados tabulares e serve como \textit{baseline} autoregressivo nos experimentos do Capítulo~\ref{cap:investigacao}.

Modelos alternativos foram considerados mas não adotados como modelos principais: o Random Forest foi avaliado no Estágio~7 como diagnóstico adicional de importância de \textit{features} via SHAP (\textit{SHapley Additive exPlanations}) \cite{lundberg2017}, mas não como modelo principal, pois o XGBoost oferece desempenho equivalente com maior controle de regularização e interpretabilidade nativa. A Regressão Logística com interações foi avaliada nos \textit{baselines} ingênuos iniciais, mas suas premissas de linearidade são inadequadas para a dinâmica não-linear de séries financeiras em janelas de 30 dias. A GRU (\textit{Gated Recurrent Unit}) foi preterida por redundância arquitetural com o BiLSTM neste contexto: ambas processam sequências com mecanismos de memória recorrente, e a literatura não evidencia ganhos sistemáticos de GRU sobre LSTM em séries financeiras de baixa frequência \cite{fischer2018}.

\section{Protocolos de avaliação em ML financeiro}

A avaliação de modelos preditivos em séries temporais financeiras apresenta desafios específicos que a distinguem de problemas padrão de \textit{machine learning}. \citeonline{bailey2014} demonstraram formalmente que a otimização de estratégias sobre um único período de \textit{backtest} inevitavelmente infla métricas de desempenho --- o fenômeno que denominaram \textit{backtest overfitting}. \citeonline{lopezdeprado2018} sistematizou o problema, propondo validação cruzada com purga e embargo como padrão mínimo. \citeonline{cawley2010} documentaram o viés de seleção de modelo, mostrando que a avaliação sobre o mesmo conjunto de dados usado para selecionar modelos produz estimativas otimistas mesmo sem vazamento explícito de dados.

A métrica ROC-AUC (\textit{Area Under the ROC Curve}) resume a capacidade de um classificador de ordenar exemplos positivos acima de negativos: equivale à probabilidade de que um par aleatório (um ``Sobe'', um ``Desce'') seja ordenado corretamente pelo escore do modelo. Vale 0,5 para um classificador aleatório e 1,0 para a separação perfeita, e não depende da escolha de um limiar de decisão. Embora amplamente utilizada, deve ser interpretada com cautela em conjuntos de teste pequenos. \citeonline{hanley1982} derivaram o erro-padrão analítico da AUC, permitindo o cálculo de intervalos de confiança e tamanhos mínimos de efeito detectáveis --- ferramenta essencial quando os conjuntos de teste têm 60--177 amostras, como neste trabalho. Como alternativa não-paramétrica ao erro-padrão analítico, este trabalho reporta intervalos de confiança via \textit{bootstrap}: o conjunto de teste é reamostrado com reposição (1.000 vezes), a AUC é recalculada em cada reamostra, e o intervalo de 95\% é lido diretamente dos percentis da distribuição resultante --- procedimento que não depende de uma fórmula fechada para a distribuição da métrica.

A validação cruzada com janela expansiva (\textit{expanding-window CV}) é a abordagem recomendada por \citeonline{lopezdeprado2018} para séries temporais financeiras: a cada \textit{fold}, o conjunto de treino expande-se cronologicamente enquanto a janela de teste avança, simulando o processo de decisão em tempo real. Em contraste com o \textit{walk-forward split} único, esta abordagem fornece múltiplas estimativas de desempenho fora da amostra, permitindo avaliar tanto o nível médio quanto a variância do desempenho ao longo do tempo.


% ============================================================================
